% ============================================================================
% WORKSHEET - Section 1.6: Applying Proportional Reasoning
% CCSS 7.RP.A.2, 7.RP.A.3
% ============================================================================
\section{Applying Proportional Reasoning}
\setcounter{wsNumber}{0}

\begin{worksheetSkills}
\begin{itemize}
    \item Solve real-world proportional problems using unit rates, equations, and scale factors
    \item Choose an efficient strategy and justify why it works
    \item Compare two proportional situations and decide which is the better value or fit
\end{itemize}
\end{worksheetSkills}

\worksheetHeader{Choose the Best Tool}

\begin{sortCut}{Sort the Problems by the Strategy You Would Use First}

Cut the cards and sort them into the strategy that feels most natural.

\bigskip
\cutCard{A car travels $210$ miles on $6$ gallons. How far on $10$ gallons?}
\cutCard{A recipe for $4$ people uses $\frac{3}{2}$ cups of rice. How much for $10$ people?}
\cutCard{On a map, $3$ cm represents $24$ km. What does $8$ cm represent?}
\cutCard{A printer makes $120$ pages in $8$ minutes. How long for $450$ pages?}
\cutCard{A concert ticket costs $\$18$. Write a rule for the total cost of $n$ tickets.}
\cutCard{A model bridge uses a scale of $1 : 200$. A beam is $6$ cm on the model. Find the real length.}

\begin{multicols}{3}
\wsSortBin[wsBlue]{Unit Rate}
\wsSortBin[wsTeal]{Equation}
\wsSortBin[wsOrange]{Scale Factor / Proportion}
\end{multicols}
\end{sortCut}
\worksheetAnswer{One strong sort is: Unit Rate - car mileage, rice recipe, and printer pages because each can be solved by finding ``per $1$'' first. Equation - concert tickets because the task asks for a rule, so $c = 18n$. Scale Factor / Proportion - the map problem and the model bridge because each compares matching measurements with a scale. Other correct sorts are possible if the reasoning is clear, but each problem must still be solved proportionally. For example, the car problem gives $210 \div 6 = 35$ miles per gallon, so $35 \times 10 = 350$ miles. The model bridge gives $6 \times 200 = 1200$ cm in real life.}

\worksheetHeader{Rate Table Decisions}

\begin{inputOutput}{Fill the Table and Use It to Solve}

\textbf{Fruit Snacks:} $5$ packs cost $\$8.75$.

\begin{center}
\begin{tabular}{|C{3cm}|C{3cm}|}
\hline
\textbf{Packs} & \textbf{Cost} \\
\hline
$5$ & $8.75$ \\
\hline
$1$ & \answerBlank[2cm] \\
\hline
$8$ & \answerBlank[2cm] \\
\hline
$12$ & \answerBlank[2cm] \\
\hline
\end{tabular}
\end{center}

Equation: \answerBlank[4cm]

\worksheetSep

If you have $\$21$, how many packs can you buy? \answerBlank[4cm]
\end{inputOutput}
\worksheetAnswer{First find the unit rate: $8.75 \div 5 = 1.75$, so one pack costs $\$1.75$. The table values are: $1$ pack costs $1.75$, $8$ packs cost $8 \times 1.75 = 14$, and $12$ packs cost $12 \times 1.75 = 21$. The equation is $c = 1.75p$. If you have $\$21$, divide by the unit price: $21 \div 1.75 = 12$, so you can buy $12$ packs.}

\worksheetHeader{Reasoning Error Patrol}

\begin{errorDetective}{Fix the Strategy Mistakes}

Amir solved the problems below.

\bigskip

Problem: $5$ shirts cost $\$42.50$. How much do $8$ shirts cost?
\errorWork{$42.50 + 8 = 50.50$}
Correct or wrong? \answerBlank[2cm] \quad Correct answer: \answerBlank[5cm]

\bigskip

Problem: $3$ cm on a map equals $24$ km. How far is $8$ cm?
\errorWork{$24 - 3 = 21$, so $8$ cm must equal $21$ km.}
Correct or wrong? \answerBlank[2cm] \quad Correct answer: \answerBlank[5cm]

\bigskip

Problem: Small paint can: $\frac{3}{4}$ gallon for $\$18$. Large paint can: $1\frac{1}{2}$ gallons for $\$33$. Which is the better deal?
\errorWork{The large can is better because $33 > 18$.}
Correct or wrong? \answerBlank[2cm] \quad Better deal: \answerBlank[5cm]
\end{errorDetective}
\worksheetAnswer{All three are wrong. (1)~The correct strategy is to find the unit price: $42.50 \div 5 = 8.50$ dollars per shirt. Then $8.50 \times 8 = 68$, so $8$ shirts cost $\$68$. (2)~The map uses a proportional scale, not subtraction. Since $24 \div 3 = 8$, each centimeter stands for $8$ km. Then $8 \times 8 = 64$ km. (3)~To compare deals, compare unit prices. Small can: $18 \div 0.75 = 24$ dollars per gallon. Large can: $33 \div 1.5 = 22$ dollars per gallon. The large can is the better deal because it costs $\$2$ less per gallon.}

\worksheetHeader{Field Trip Budget}

\begin{mathScenario}{Plan the Field Trip Snacks}

The seventh-grade team is buying granola bars for a field trip.

\bigskip
\scenarioItem{Store A}{$6$ bars for $\$5.10$}
\scenarioItem{Store B}{$10$ bars for $\$8.20$}
\scenarioItem{Needed for the trip}{$40$ bars}

\bigskip

What is the unit price at Store A? \answerBlank[4cm]

\bigskip

What is the unit price at Store B? \answerBlank[4cm]

\bigskip

Which store should the team choose for $40$ bars, and how much money will that store cost? \answerBlank[8cm]
\end{mathScenario}
\worksheetAnswer{(1)~Store A unit price: $5.10 \div 6 = 0.85$, so each bar costs $\$0.85$. (2)~Store B unit price: $8.20 \div 10 = 0.82$, so each bar costs $\$0.82$. (3)~Store B is cheaper. For $40$ bars, multiply the unit price: $40 \times 0.82 = 32.80$. Store A would cost $40 \times 0.85 = 34.00$. The team should choose Store B and spend $\$32.80$.}

\worksheetHeader{Explain Your Strategy Choice}

\begin{explainIt}{Why Was That Strategy Efficient?}

Think about the problem ``A recipe for $4$ people uses $\frac{3}{2}$ cups of rice. How much for $10$ people?''
Explain why finding the amount for one person first is a good strategy.

\writeLines{5}

\bigskip

A classmate solved a proportion problem with an equation, but you solved it with a unit rate. Explain why two different strategies can still be correct.

\writeLines{6}
\end{explainIt}
\worksheetAnswer{(1)~Finding the amount for one person first is efficient because it creates a unit rate that can be scaled to any group size. For the rice problem, $\frac{3}{2} \div 4 = \frac{3}{8}$ cup per person, and then multiplying by $10$ gives $\frac{30}{8} = 3\frac{3}{4}$ cups. Once the unit rate is known, it is easy to build up or compare. (2)~Different strategies can still be correct because they are based on the same proportional relationship. A unit rate, an equation like $y = kx$, and a proportion all describe the same constant ratio. If each method is set up correctly, they lead to the same answer.}

\worksheetDone